\subsection{An exact optimized cost}
Throughout this section put
\begin{equation}\label{adaptive:range}
 \frac1{10}\le a\le\frac16,\qquad 2a\le b\le\frac12,
 \qquad 0\le\theta\le1.
\end{equation}
The parameter \(\theta\) interpolates between a fixed head interval and an
interval determined by the actual tail minimum. Define
\begin{align*}
 k&=2b-a,&q&=\frac b{k},&
 h_\theta&=\frac{1+q}{2}-\frac{\theta a}{4k},&
 r_\theta&=\frac\theta{4b},\\
 L_0&=\frac{1-2a}{4},&\lambda&=\frac{1-2b}{4b},&
 A_0&=\frac{1-a-3b}{3},&\gamma&=\frac{6b-1-2a}{3a}.
\end{align*}
Here \(\gamma\ge0\), and \(\gamma+\lambda>0\), including the corner
\(a=1/10,b=1/5\), where \(\gamma=0\). Set
\begin{align}
 m_\theta&=\frac{a h_\theta}{1-a r_\theta},&
 E_\theta&=A_0+\gamma m_\theta,&
 J_E&=\frac{\gamma L_0+\lambda A_0}{\gamma+\lambda},
 \label{adaptive:early}\\
 F&=\frac{(1-a)(b-a)}{4(2b-a)(1-b)},&
 Z&=\min\left\{\frac{1-q}{3},F\right\}.
 \label{adaptive:zero}
\end{align}
For the two late routes, indexed by \(v\) (variation) and \(\ell\) (length), put
\begin{align*}
 A_v&=\frac{1-a}{3}-\frac{(1-b)h_\theta}{2},&
 \sigma_v&=\frac{1-a}{6a}-\frac{(1-b)r_\theta}{2},\\
 A_\ell&=\frac{1-a}{3}-h_\theta,&
 \sigma_\ell&=\frac{a+2}{3a}-r_\theta,\\
 P_i&=A_i+a\sigma_i,&
 J_i&=\frac{\sigma_iL_0+\lambda A_i}{\sigma_i+\lambda}
 \quad(i=v,\ell).
\end{align*}
Both \(\sigma_i\) are strictly positive: for example
\(\sigma_v\ge(5-5a)/(48a)>0\), using \(b\ge2a\).
Also \(1-a r_\theta\ge7/8\). Thus all denominators are positive.
Our new cost and its compact optimization are
\begin{align}
 C_\theta(a,b)&=\max\left\{Z,\ \min(E_\theta,J_E),\
                  \min(P_v,J_v,P_\ell,J_\ell)\right\},
                   \label{adaptive:cost}\\
 C_{\rm A}(a,b)&=\min_{0\le\theta\le1}C_\theta(a,b).
                   \label{adaptive:optimized}
\end{align}
These are exact real expressions. The last minimum exists by continuity on
the compact interval; it is not an assumption about a numerical optimizer.

\begin{lemma}[Normalized adaptive profile bound]\label{adaptive:profile}
Under \eqref{adaptive:range}, let \(g:[0,N]\to\R\) be 1-Lipschitz and
\(ax\le g(x)\le bx\). For \(15N/16\le n_0<N\), there is an admissible
chain from \(n_0\) to zero of cost at most \(NC_\theta(a,b)\), and hence
one of cost at most \(NC_{\rm A}(a,b)\). The number of edges is
\(O(1+\log(N/(N-n_0)))\), uniformly in \(\theta\). Grid and approximate
barrier errors are \(O(KT+K\epsilon N)\), as in Lemma~\ref{minimum:profile}.
\end{lemma}

\subsection{Endpoint normalization and free continuation}
Normalize \(N=1\). Replace \(g(x)\) by
\(\widetilde g(x)=\min\{g(x),1+a-x\}\). This preserves the barriers
and the Lipschitz constant, and sets \(\widetilde g(1)=a\). Every edge
cost increases. Indeed, if its left endpoint \(v\) is unchanged, the
interval minimum can only decrease. Otherwise
\[
 c_{\widetilde g}(v,n)\ge(1+a-v)-(1+a-n)=n-v\ge c_g(v,n).
\]
It is therefore enough to construct chains for a profile with \(g(1)=a\).

Let \(t\) minimize \(g\) on \([q,1]\), and write \(m=g(t)\). Then
\begin{equation}\label{adaptive:minimum-bounds}
 aq\le at\le m\le a,\qquad m\le bq.
\end{equation}
Every endpoint below \(q\) has a free continuation, by
\(g(2n-1)\le b(2n-1)\le an\le g(n)\) for \(1/2<n\le q\),
and the free jump to zero below \(1/2\). The free region extends to
\[
 Q_m=\frac{1+m/b}{2},
\]
because \(g(2n-1)\le m\le g(n)\) when \(q<n\le Q_m\).
Set
\[
 Q_\theta=(1-\theta)q+\theta Q_m,
 \qquad H_\theta=\frac{1+Q_\theta}{2}=h_\theta+r_\theta m.
\]
Thus every endpoint below \(Q_\theta\) has a free continuation,
\(q\le Q_\theta\le Q_m\), and
\(H_\theta\le(3+q)/4\le11/12<n_0\).

The affine majorant \(z/2+(b-1/2)m/b\), applied above \(Q_m\),
has nonpositive potential at \(Q_m\). The potential construction
\eqref{minimum:potential} therefore supplies, independently of all later
cases, a chain with cost
\begin{equation}\label{adaptive:floor}
 L(m)=L_0-\lambda m.
\end{equation}

\subsection{Early minima and the supporting line}
If \(t\le Q_\theta\), the cone chain to \(t\) followed by free
continuation costs at most \((1-t+m-a)/3\). If
\(Q_\theta<t\le H_\theta\), jump after that cone to a minimum on
\([2t-1,Q_\theta]\), and then continue freely. The cost is at most
\begin{equation}\label{adaptive:early-route}
 \frac{1-t+m-a}{3}+\bigl(b(2t-1)-m\bigr)_+.
\end{equation}

When the terminal contribution is zero, the cone bound is at most
\((1-q)/3\). A second construction gives \(F\). On \([t,1]\),
\(g(z)\le\min\{bz,m+z-t\}\). These two lines meet at
\(z_*=(t-m)/(1-b)\). The affine function
\[
 H(z)=\frac z2+\left(b-\frac12\right)z_*
\]
majorizes their minimum. Its potential bounds the cost down to
\((1+t)/2\), followed by a free jump to \(t\), by
\[
 \frac{H(1)-a}{2}
 =\frac{1/2-a-\kappa(t-m)}2,
 \qquad\kappa=\frac{1/2-b}{1-b}\ge0.
\]
Here the construction can stop with positive remaining potential: its
discarded value is nonnegative. If the potential reaches zero first,
use zero-cost minimum steps until the terminal jump is admissible. All
left endpoints where the majorant is used are at least \(t\).
Since \(t-m\ge q-a\), the displayed expression is at most \(F\), by
direct substitution. The zero-terminal case thus has cost at most \(Z\).

If the terminal contribution in \eqref{adaptive:early-route} is positive,
the bound becomes
\[
 A(t,m)=\frac{1-a-3b+(6b-1)t-2m}{3}.
\]
Using \(t\le m/a\) gives the increasing line \(A_0+\gamma m\).
Using \(t\le H_\theta\) gives a decreasing line, whose slope is
\(((6b-1)r_\theta-2)/3<0\). They intersect at \(m_\theta\),
with common value \(E_\theta\). Comparing the increasing line with
the independent floor \eqref{adaptive:floor} also gives \(J_E\).
Hence this case has cost at most \(\min(E_\theta,J_E)\).

\subsection{Late minima: reflection and two terminal routes}
Suppose \(t>H_\theta\), and let \(u\) minimize \(g\) on
\([Q_\theta,H_\theta]\). Put \(w=g(u)\), \(z=b(2u-1)\).
Then \(w\ge m\), \(2u-1\le Q_\theta\), and
\(m\ge at>aH_\theta\), so \(m>m_\theta\).

If \(w\le z\), use the running-minimum reflection from the preceding
section: set \(r=0\) before \(u\),
\(r(x)=w-\min_{[u,x]}g\) after \(u\), and \(f=g+r\).
It is 1-Lipschitz, has minimum \(w\) on \([Q_\theta,1]\), and
\(f(1)=a+w-m\). The edge inequality
\(c_g(v,n)\le c_f(v,n)+r(n)-r(v)\) transfers a cone chain to \(u\),
its terminal jump, and free continuation, with total cost at most
\[
 \frac{1-u-a-2m}{3}+\max(w,z)
 =\frac{1-a-3b+(6b-1)u-2m}{3}.
\]
Since \(u\le H_\theta<m/a\), both previous affine bounds apply;
the decreasing one is evaluated to the right of \(m_\theta\).
The bound is again \(\min(E_\theta,J_E)\).

If \(w\ge z\), the interval below \(H_\theta\) is free: jump first
to \(u\), then to a minimum on \([2u-1,Q_\theta]\), then continue
freely. Above it, use the cone to \(t\) and an admissible chain from
\(t\) to \(H_\theta\). Bounding this middle part by negative variation
or by its length gives respectively
\[
 \frac{1-a}{3}-\frac{(1-b)H_\theta}{2}+\frac{t-m}{6},
 \qquad
 \frac{1-a}{3}-H_\theta+\frac{2t+m}{3}.
\]
The relation \(t\le m/a\) bounds these by
\(A_i+\sigma_i m\), \(i=v,\ell\). If \(t>n_0\), omit the cone
and restrict the middle variation or length to \([H_\theta,n_0]\);
the omitted cone allowance is nonnegative, so the bounds remain valid.
Since \(m\le a\), each increasing line is at most \(P_i\).
Intersecting it with the decreasing floor gives \(J_i\). The late
free-terminal case therefore has cost at most
\(\min(P_v,J_v,P_\ell,J_\ell)\), completing the cost proof.

There are a bounded number of chain portions. For clarity, compression
does not require that their edges be free: for \(l<m<n\),
\[
 c_g(l,n)\le c_g(l,m)+c_g(m,n).
\]
This follows by comparing the two interval minima and using
\(\min_{[l,m]}g\le g(m)\). Merge any adjacent pair whose union is
admissible. In a chain with no further such merge, complementary depth
more than doubles every two edges, giving the claimed logarithmic count.
For piecewise affine profiles the constructions terminate; approximate
a general Lipschitz profile by piecewise affine interpolation and pass to
a convergent subsequence of uniformly bounded endpoint lists. Costs,
minima and admissibility pass to the limit. Rounding down and clamping
the barriers give the stated errors exactly as before. All bounds are
uniform in \(\theta\); this proves Lemma~\ref{adaptive:profile}.

\subsection{The unique root and the dimension curve}
\begin{lemma}\label{adaptive:root}
The function \(C_{\rm A}\) is continuous and strictly increasing in \(b\)
for fixed \(a\) in \eqref{adaptive:range}. There is a unique continuous
\(\beta(a)\in[2a,1/2]\) satisfying \(C_{\rm A}(a,\beta(a))=a\).
Its endpoint values are \(\beta(1/10)=1/5\), \(\beta(1/6)=1/2\),
and \(\beta(a)>2a\) for \(1/10<a\le1/6\).
\end{lemma}
\begin{proof}
All expressions are jointly continuous with positive denominators;
minimizing over a fixed compact interval preserves continuity.
We check strict increase of every ingredient in \eqref{adaptive:cost}.
Both terms defining \(Z\) increase strictly. Write
\[
 E_\theta=\frac{b-a}{3(2b-a)}T,\qquad
 T=\frac{12b^2+(2-8a)b-\theta a(1-a)}{4b-\theta a}.
\]
The first factor increases strictly. We have \(T>0\), and its derivative
has positive numerator
\(48b^2-24ab\theta+2a\theta+4a^2\theta\).
Thus \(E_\theta\) increases strictly. For \(J_E\), its derivative is
a positive factor times \(60b^2-12b+1-4a\). Substitute
\(b=2a+x\), \(x\ge0\), to obtain
\[
 240a^2-28a+1+(240a-12)x+60x^2>0.
\]
Indeed, the constant term is increasing for \(a\ge1/10\) and is
\(3/5\) at \(a=1/10\). Thus \(J_E\) increases strictly.

For the late bounds, regard \(H_\theta(m)=h_\theta+r_\theta m\)
as a function of \(b\) at fixed \(a,\theta,m>0\). Then
\[
 \partial_b H_\theta(m)
 =-\frac{a(1-\theta)}{2(2b-a)^2}-\frac{\theta m}{4b^2}<0.
\]
Both route lines \(A_i+\sigma_i m\) therefore increase strictly in
\(b\), as does \(L(m)\), since \(\partial_b L(m)=m/(4b^2)>0\).
At \(m=aq\), the differences between the floor and the two route lines are
\[
 \frac{(b-a)(5-2a-9b)}{12(2b-a)}>0,
 \qquad \frac{(4-a)(b-a)}{6(2b-a)}>0.
\]
Their intersections are consequently at positive \(m\). To compare
intersection values at two upper slopes, evaluate the two new lines at
the old intersection: both values have increased strictly. The maximum
of the minimum of a decreasing and an increasing affine line is their
intersection value, so \(J_i\) increases strictly. Evaluating the route
lines at \(m=a\) similarly proves this for \(P_i\).
Finite minima and maxima preserve strict increase. Finally, choosing a
minimizing \(\theta\) at the larger upper slope proves strict increase
of \(C_{\rm A}\) itself.

For the lower endpoint, use \(\theta=1,b=2a\). Direct substitution
gives a zero-terminal bound
\((1-a)/(12(1-2a))<a\), an early bound at most
\((1+11a)/21\), and a late bound at most \((1+6a)/16\).
The latter two are at most \(a\), strictly so for \(a>1/10\).
At \(a=1/10,b=1/5\), \(\gamma=0\) forces the early bound to be
exactly \(a\) for every \(\theta\), establishing equality there.

At \(b=1/2\), \(\lambda=0\), so all three intersection values equal
\(L_0\). The other early bound is at least \(E_0=4L_0/3\).
The zero-terminal bound is at most \(F=L_0\), while
\(P_\ell\ge L_0\) and \(P_v\ge L_0\).
Consequently \(C_\theta(a,1/2)=L_0\) for every \(\theta\).
Now \(L_0\ge a\), with equality only at \(a=1/6\).
The intermediate value theorem and strict increase give the unique root.
Continuity of the root follows by taking convergent subsequences in the
compact interval and using uniqueness, including the two endpoints.
\end{proof}

This gives the new curve \(B_{\rm A}\) stated at the beginning of the
manuscript. It dominates \(B_{\rm M}\): on the latter's nontrivial middle
range, \(\theta=0\) gives the same early bounds as \(V\), no larger
late bounds, and \(Z\le(1-q)/3\le\min(E_2,J_E)\). Thus
\(C_{\rm A}\le C_0\le V\). Below that range the previous curve is
\(2d-1\); above it the two displayed curves agree.

An exact new example, requiring no optimization, is
\[
 a=\frac{11}{100},\quad b=\frac9{40},\quad\theta=1,
 \qquad C_1(a,b)=\frac{644}{5925},\qquad
 a-C_1(a,b)=\frac{31}{23700}>0.
\]
The fully interpolated construction also covers
\(a=7/50,b=11/35,\theta=1/3\), with cost
\(452437/3234750\) and margin \(214/1617375\).
These are rational checks of strict inequalities, not numerical evidence
substituting for the profile or analytic proof.
